%%%%%%%%%%%%%%%%%%%%%%%%%%%%
% PAQUETES
%%%%%%%%%%%%%%%%%%%%%%%%%%%%

\usepackage[utf8]{inputenc}                % Codificación UTF-8 para caracteres especiales
\usepackage[T1]{fontenc}                   % Mejora la salida de texto con acentos
% Babel no está disponible en esta instalación de TeX; se desactiva para que el documento compile.
% Si deseas activar el idioma español, instala el paquete correspondiente del sistema TeX y descomenta la línea:
% \usepackage[spanish]{babel}
\usepackage{graphicx}                      % Manejo avanzado de imágenes [\includegraphics]
\usepackage{float}                         % Posicionamiento preciso de figuras/tablas [H]
\usepackage{bookmark}                      % Bookmarks interactivos en PDF (mejorado)
\usepackage[makeroom]{cancel}              % Cancelación de términos en ecuaciones [\cancel]
\usepackage{hyperref}                      % Hipervínculos interactivos
\usepackage{nameref}                       % Referencias por nombre de sección
\usepackage{theoremref}                    % Referencias a teoremas mejoradas
\hypersetup{
    pdftitle={Assignment},                 % Metadatos del PDF
    colorlinks=true,                       % Enlaces en color (no recuadros)
    linkcolor=cyan!90,                     % Color enlaces internos
    urlcolor=blue,                         % Color URLs externas
    bookmarksnumbered=true,                % Numeración automática de bookmarks
    bookmarksopen=true                     % Bookmarks expandidos por defecto
}

\usepackage{amsmath, amsfonts, amsthm, amssymb, mathtools} % Base para matemáticas
\usepackage{xfrac}                         % Fracciones diagonales [\sfrac]
\usepackage{tikz-cd}                       % Diagramas conmutativos
\usepackage{tikzsymbols}                   % Símbolos especiales (emojis, misc.)
\usepackage{siunitx}                       % Alineación numérica en tablas [S]
\usepackage{multicol}                      % Columnas múltiples en texto/matemáticas

\usepackage{booktabs}                      % Reglas profesionales para tablas [\toprule]
\usepackage{array}                         % Control avanzado de columnas en tablas
\usepackage[ruled,vlined,linesnumbered]{algorithm2e} % Escritura de algoritmos
\usepackage[most,many,breakable]{tcolorbox} % Cajas decoradas con estilo
\usepackage{xcolor}                         % Gestión avanzada de colores
\usepackage{varwidth}                       % Cajas de ancho variable
\usepackage{enumitem}                       % Personalización de listas [nosep]
\usepackage{etoolbox}                      % Utilidades para programación LaTeX
\usepackage{import}                        % Importación modular de archivos
\usepackage{xifthen}                       % Condicionales en documentos
\usepackage{pdfpages}                     % Inclusión de páginas PDF externas
\usepackage{transparent}                  % Manejo de transparencias
% `babelbib` requiere soporte de idioma de babel, que no está instalado en este sistema.
% Se desactiva para que el documento compile sin errores.
% \usepackage[fixlanguage]{babelbib}
% \bibliographystyle{babunsrt}
\usepackage{titlesec}
% Personalización de títulos
% Centrar secciones automáticamente
\titleformat{\section}
  {\normalfont\Large\bfseries\centering}{\thesection}{1em}{}
\titleformat{\subsection}
  {\normalfont\large\bfseries\centering}{\thesubsection}{1em}{}

\usepackage{comment}  
\usepackage{epsfig, graphicx} % Inclusión de imágenes y gráficos
\usepackage[left=2cm, right=2cm, top=1.8cm, bottom=2.3cm]{geometry} % Define los márgenes del documento
\usepackage{fancyhdr}         % Para personalizar encabezados y pies de página
\usepackage{lastpage}         % Permite referenciar el número total de páginas
\usepackage{lipsum}           % Genera texto aleatorio para pruebas
\usepackage{pgfplots, mathtools}


%%%%%%%%%%%%%%%%%%%%%%%%%%%%
% COLORES
%%%%%%%%%%%%%%%%%%%%%%%%%%%%

\definecolor{mygr}{HTML}{2C3338}
\definecolor{mauve}{rgb}{0.58,0,0.82}  % Color morado

%%%%%%%%%%%%%%%%%%%%%%%%%%%%
% CAJAS
%%%%%%%%%%%%%%%%%%%%%%%%%%%%

% --------------------------
% Caja de Introducción/Indicaciones
% --------------------------
\newtcolorbox{intro}[1][]{%
    enhanced,                               % Habilitar mejoras gráficas
    before skip=2mm, after skip=2mm,        % Espaciado vertical
    colback=gray!5,                         % Color de fondo claro
    colframe=black!80!black,                % Borde oscuro
    boxrule=0.5mm,                          % Grosor del borde
    attach boxed title to top left={         % Posición del título
        xshift=1cm,                         % Desplazamiento horizontal
        yshift*=1mm-\tcboxedtitleheight      % Ajuste vertical preciso
    },
    varwidth boxed title*=-3cm,             % Ancho adaptable del título
    boxed title style={                     % Estilo personalizado del título
        frame code={                        % Diseño del marco del título
            \path[fill=tcbcolback]
            ([yshift=-1mm,xshift=-1mm]frame.north west)
            arc[start angle=0,end angle=180,radius=1mm]
            ([yshift=-1mm,xshift=1mm]frame.north east)
            arc[start angle=180,end angle=0,radius=1mm];
            \path[left color=tcbcolback!60!black,
                  right color=tcbcolback!60!black,
                  middle color=tcbcolback!80!black]
            ([xshift=-2mm]frame.north west) -- ([xshift=2mm]frame.north east)
            [rounded corners=1mm]-- ([xshift=1mm,yshift=-1mm]frame.north east)
            -- (frame.south east) -- (frame.south west)
            -- ([xshift=-1mm,yshift=-1mm]frame.north west)
            [sharp corners]-- cycle;
        }, 
        interior engine=empty,              % Sin relleno adicional
    },
    fonttitle=\bfseries,                    % Título en negrita
    title={Indicaciones},                   % Texto predeterminado del título
    #1                                      % Permite personalización adicional
}

% --------------------------
% Caja de Bibliografía (similar a intro con título diferente)
% --------------------------
\newtcolorbox{biblio}[1][]{enhanced,
    before skip=2mm, after skip=2mm, colback=gray!5, colframe=black!80!black, boxrule=0.5mm,
    attach boxed title to top left={xshift=1cm, yshift*=1mm-\tcboxedtitleheight}, varwidth boxed title*=-3cm,
    boxed title style={frame code={
            \path[fill=tcbcolback]
            ([yshift=-1mm,xshift=-1mm]frame.north west)
            arc[start angle=0,end angle=180,radius=1mm]
            ([yshift=-1mm,xshift=1mm]frame.north east)
            arc[start angle=180,end angle=0,radius=1mm];
            \path[left color=tcbcolback!60!black,right color=tcbcolback!60!black,
                middle color=tcbcolback!80!black]
            ([xshift=-2mm]frame.north west) -- ([xshift=2mm]frame.north east)
            [rounded corners=1mm]-- ([xshift=1mm,yshift=-1mm]frame.north east)
            -- (frame.south east) -- (frame.south west)
            -- ([xshift=-1mm,yshift=-1mm]frame.north west)
            [sharp corners]-- cycle;
        }, interior engine=empty,
    },
    fonttitle=\bfseries,
    title={Bibliografía}, #1
}

% --------------------------
% Caja de Inciso
% --------------------------
\makeatletter
\newtcolorbox{inciso}[2][]{%
    enhanced,
    breakable,
    colback=white,
    colframe=mygr,
    attach boxed title to top left={
        yshift*=-\tcboxedtitleheight
    },
    fonttitle=\bfseries,
    title={#2},
    boxed title size=title,
    boxed title style={
        sharp corners,
        rounded corners=northwest,
        colback=tcbcolframe,
        boxrule=0pt,
    },
    underlay boxed title={
        \path[fill=tcbcolframe] (title.south west)--(title.south east)
        to[out=0, in=180] ([xshift=5mm]title.east)--
        (title.center-|frame.east)
        [rounded corners=\kvtcb@arc] |-
        (frame.north) -| cycle;
    },
    #1
}
\makeatother

% --------------------------
% Caja de Notas
% --------------------------
\usetikzlibrary{arrows,calc,shadows.blur}   % Requisitos para efectos gráficos
\tcbuselibrary{skins}                       % Habilita skins avanzados

\newtcolorbox{notas}[1][]{%
    enhanced jigsaw,                        % Skin mejorado para efectos
    colback=gray!20!white,                  % Fondo gris claro
    colframe=gray!80!black,                 % Borde gris oscuro
    size=small,                             % Tamaño compacto
    boxrule=1pt,                            % Grosor del borde
    title=\textbf{Notas | Observaciones},   % Título doble
    halign title=flush center,              % Alineación central del título
    coltitle=black,                         % Color del texto del título
    breakable,                              % Permite dividir en páginas
    drop shadow=black!50!white,             % Sombra exterior
    attach boxed title to top left={         % Posición del título
        xshift=1cm,
        yshift=-\tcboxedtitleheight/2,
        yshifttext=-\tcboxedtitleheight/2
    },
    minipage boxed title=5cm,               % Ancho máximo del título
    boxed title style={                     % Estilo personalizado del título
        colback=white,                      % Fondo blanco
        size=fbox,                         % Tamaño tipo marco
        boxrule=1pt,                        % Grosor del borde
        boxsep=2pt,                         % Espaciado interno
        underlay={                          % Elementos gráficos bajo el título
            \coordinate (dotA) at ($(interior.west) + (-0.5pt,0)$);
            \coordinate (dotB) at ($(interior.east) + (0.5pt,0)$);
            \begin{scope}
                \clip (interior.north west) rectangle ([xshift=3ex]interior.east);
                \filldraw [white, blur shadow={shadow opacity=60, shadow yshift=-.75ex}, 
                          rounded corners=2pt] (interior.north west) rectangle (interior.south east);
            \end{scope}
            \begin{scope}[gray!80!black]
                \fill (dotA) circle (2pt);  % Punto decorativo izquierdo
                \fill (dotB) circle (2pt);  % Punto decorativo derecho
            \end{scope}
        },
    },
    #1                                      % Personalización adicional
}

% --------------------------
% Caja de Demostración
% --------------------------
\newtcbtheorem[]{demostracion}{Demostracion}%
{ enhanced,
  colback       = black!5,               % Color de fondo del recuadro
  colbacktitle  = black!5,               % Color de fondo del título
  coltitle      = black,                  % Color del texto del título
  boxrule       = 0pt,                    % Ancho del borde del recuadro
  frame hidden,                           % Oculta el marco del recuadro
  borderline west = {0.5mm}{0.0mm}{black},% Línea decorativa en el lateral izquierdo
  fonttitle     = \bfseries\sffamily,      % Fuente y estilo del título
  breakable,                              % Permite que el recuadro se divida en varias páginas
  before skip   = 3ex,                    % Espacio antes del entorno
  after skip    = 3ex                     % Espacio después del entorno
}{problem}


%%%%%%%%%%%%%%%%%%%%%%%%%%%%
% CONFIGURACIONES
%%%%%%%%%%%%%%%%%%%%%%%%%%%%

% Fallback para que el documento compile aunque babel no esté instalado.
\providecommand{\spanishdecimal}[1]{ }
% Ajusta el espaciado antes y después de los títulos de capítulos y secciones.
\titlespacing*{\chapter}{0cm}{-2.0cm}{0.50cm} % {izquierda}{espacio antes}{espacio después}
\titlespacing*{\section}{0cm}{0.50cm}{0.25cm}

% Configura la indentación y separación entre párrafos.
\setlength{\parindent}{0pt}  % Sin indentación en los párrafos
\setlength{\parskip}{1ex}    % Espacio extra entre párrafos


% Configuración de encabezados y pies con fancyhdr.
\pagestyle{fancy}
\fancyhf{}  % Limpia los encabezados y pies predeterminados
\rfoot{%
  \includegraphics[height=1cm]{Logo_FC.png} \hfill % Inserta un logo a la derecha
  \textit{Página \thepage\ de \pageref{LastPage}}
}

\usetikzlibrary{arrows, backgrounds, shapes.arrows, positioning} % Librerías de TikZ para gráficos vectoriales
\pgfplotsset{compat=newest}   % Asegura la compatibilidad con la versión más reciente de pgfplots

%%%%%%%%%%%%%%%%%%%%%%%%%%%%
% DEFINICIÓN DE ENTORNOS Y COMANDOS
%%%%%%%%%%%%%%%%%%%%%%%%%%%%

% definición del operador para \left( \right)
\newcommand{\paren}[1]{\left({#1}\right)}

% definición del operador para derivada
\newcommand{\dd}[2]{\dfrac{\mathrm{d}#1}{\mathrm{d}#2}}

% definición del operador para la función techo
\newcommand{\techo}[1]{\lceil {#1} \rceil}

% Redefine la numeración de listas para usar números (permite contadores más grandes que el sistema alfabético).
\renewcommand{\theenumi}{\arabic{enumi})}

% Comandos personalizados para mejorar la notación matemática.
\let\emptyset\varnothing         % Utiliza el símbolo \varnothing para el conjunto vacío
\newcommand{\lto}{\mathbin{\to}}   % Define un comando para la flecha "→"

% Redefine los operadores de parte real e imaginaria para evitar conflictos.
\let\Re\relax
\DeclareMathOperator{\Re}{Re}
\let\Im\relax
\DeclareMathOperator{\Im}{Im}

% Definición de entornos para teoremas, corolarios, definiciones y proposiciones.
\newtheorem{teorema}{Teorema}
\newtheorem{corolario}{Corolario}
\newtheorem{definición}{Definición}
\newtheorem{proposición}{Proposición}